% Options for packages loaded elsewhere
\PassOptionsToPackage{unicode}{hyperref}
\PassOptionsToPackage{hyphens}{url}
%
\documentclass[
  10pt,
]{article}
\usepackage{amsmath,amssymb}
\usepackage{iftex}
\ifPDFTeX
  \usepackage[T1]{fontenc}
  \usepackage[utf8]{inputenc}
  \usepackage{textcomp} % provide euro and other symbols
\else % if luatex or xetex
  \usepackage{unicode-math} % this also loads fontspec
  \defaultfontfeatures{Scale=MatchLowercase}
  \defaultfontfeatures[\rmfamily]{Ligatures=TeX,Scale=1}
\fi
\usepackage{lmodern}
\ifPDFTeX\else
  % xetex/luatex font selection
\fi
% Use upquote if available, for straight quotes in verbatim environments
\IfFileExists{upquote.sty}{\usepackage{upquote}}{}
\IfFileExists{microtype.sty}{% use microtype if available
  \usepackage[]{microtype}
  \UseMicrotypeSet[protrusion]{basicmath} % disable protrusion for tt fonts
}{}
\makeatletter
\@ifundefined{KOMAClassName}{% if non-KOMA class
  \IfFileExists{parskip.sty}{%
    \usepackage{parskip}
  }{% else
    \setlength{\parindent}{0pt}
    \setlength{\parskip}{6pt plus 2pt minus 1pt}}
}{% if KOMA class
  \KOMAoptions{parskip=half}}
\makeatother
\usepackage{xcolor}
\usepackage[margin=25mm]{geometry}
\usepackage{longtable,booktabs,array}
\usepackage{calc} % for calculating minipage widths
% Correct order of tables after \paragraph or \subparagraph
\usepackage{etoolbox}
\makeatletter
\patchcmd\longtable{\par}{\if@noskipsec\mbox{}\fi\par}{}{}
\makeatother
% Allow footnotes in longtable head/foot
\IfFileExists{footnotehyper.sty}{\usepackage{footnotehyper}}{\usepackage{footnote}}
\makesavenoteenv{longtable}
\setlength{\emergencystretch}{3em} % prevent overfull lines
\providecommand{\tightlist}{%
  \setlength{\itemsep}{0pt}\setlength{\parskip}{0pt}}
\setcounter{secnumdepth}{-\maxdimen} % remove section numbering
\usepackage{amsmath,amssymb,mathtools,fvextra,xurl,microtype,needspace,longtable,booktabs,array}
\usepackage{scrlayer-scrpage}
\clearpairofpagestyles
\ihead{Split-Zero / Benzel P7}\ohead{Cumulative research record}\cfoot{\pagemark}
\setkomafont{pageheadfoot}{\normalfont\small}
\setkomafont{disposition}{\normalfont\bfseries}
\setlength{\emergencystretch}{3em}
\DefineVerbatimEnvironment{Highlighting}{Verbatim}{breaklines,breakanywhere,commandchars=\\\{\},fontsize=\footnotesize}
\RecustomVerbatimEnvironment{verbatim}{Verbatim}{breaklines,breakanywhere,fontsize=\footnotesize}
\fvset{breaklines=true,breakanywhere=true,fontsize=\footnotesize}
\AtBeginEnvironment{longtable}{\footnotesize}
\ifLuaTeX
  \usepackage{selnolig}  % disable illegal ligatures
\fi
\usepackage{bookmark}
\IfFileExists{xurl.sty}{\usepackage{xurl}}{} % add URL line breaks if available
\urlstyle{same}
\hypersetup{
  hidelinks,
  pdfcreator={LaTeX via pandoc}}

\author{}
\date{}

\begin{document}

\section{An asymmetric positive completion for the unbounded deletion
family q = 3k +
1}\label{an-asymmetric-positive-completion-for-the-unbounded-deletion-family-q-3k-1}

\textbf{Benzel P7 continuation, 17 September 2026.} This note gives an
original-coordinate, certificate-assisted construction for every integer
\(k\ge4\), \(m\ge3k+3\), and \(d\ge m\). It does not declare the full
Propp Problem 7 resolved. Independent specialist review, a Lean
reconstruction, and a novelty determination have not been completed. The
exact turn-8 source remains unchanged alongside this continuation.

\subsection{1. The completed scope and original
objects}\label{the-completed-scope-and-original-objects}

Put \(t_n=n(n-1)/2\) and \[
q=3k+1,\qquad \Delta=t_m-q,\qquad h=t_d-\Delta.
\] For all the parameters above, the construction in this note tiles \[
W_h(d)=V(d+3h,2d+3h)
\] with exactly \(\Delta\) original right stones and \(3h(h+d)\)
original bones. Reflection supplies the reversed parameter pair. Thus
\(q=13,16,19,\ldots\) is treated without an upper bound on \(q\). The
cases \(k=1,2,3\) are not included in this new theorem; earlier
fixed-deletion constructions retain their own scope.

The original axial cell \((x,y)\) has differences \[
(y-x,\;1-x-2y,\;2x+y-1).
\] The benzel \(V(a,b)\) requires all three differences to lie in
\([1-a,b-1]\). Its barycentric cell is \((x,y,1-x-y)\), with projection
onto the first two coordinates as inverse. The original region-parameter
inverse is \(d=b-a\), \(h=(2a-b)/3\) on this integer lane.

The permitted tiles and their original cell offsets are

\begin{longtable}[]{@{}ll@{}}
\toprule\noalign{}
Tile & Offsets \\
\midrule\noalign{}
\endhead
\bottomrule\noalign{}
\endlastfoot
\(R\) & \((0,0),(1,0),(0,1)\) \\
\(H\) & \((0,0),(1,0),(2,0)\) \\
\(V\) & \((0,0),(0,1),(0,2)\) \\
\(D\) & \((0,0),(1,-1),(2,-2)\) \\
\end{longtable}

The incidence map \(\partial\) sends each placed tile to the sum of its
three original cell generators. Rotation and reflection are \[
\rho(x,y)=(y,1-x-y),\qquad F(x,y)=(x,1-x-y).
\] Their inverses are \(\rho^2\) and \(F\). A tiling is the nonnegative
integral fibre of \(\partial\) over the original all-ones region vector,
not an arbitrary signed chain with the same image.

\subsection{2. The parent packing and the exact old
support}\label{the-parent-packing-and-the-exact-old-support}

We use the complete stable-packing proof retained in
\path|../turn8/sources/turn6-PROOF.md|. Its actual incidence identity is
\[
\partial\Gamma_\Delta(d,h)+\mathbf1_{\Omega_\Delta}
 =\mathbf1_{W_h(d)}.
\] Here \(\Omega_\Delta\) consists of the first \(\Delta\) original rank
orbits. The rank representative is \[
p_{r,s}=(r-s,-s),\quad 0\le s<r,\quad
\operatorname{rank}(p_{r,s})=t_r+s+1,
\] with its two rotations. Selecting the position of the minimum
difference, its negative value \(r\), and \(s\) gives the inverse map on
the original cells.

The common old bone family \(A\) is indexed by \(p\in\{0,1,2\}\),
\(1\le y\le m+q\), and \(0\le\ell<\min(y,q)\). Its actual generator is
\[
a_{p,y,\ell}=F\rho^p H(y-r_y-2-3\ell,y),
\qquad
r_y=\begin{cases}m-1&y\le q,\\m&y>q.\end{cases}
\] These are disjoint original bones outside \(\Omega_\Delta\). This
entire finite family is retained; the new proof does not claim that its
release count is minimal.

\subsubsection{2.1 The source-to-tail receiving
maps}\label{the-source-to-tail-receiving-maps}

At \(d=m\), \(h=q\). Prefix rows use index \(j=y-1-\ell\). In a tail row
\(q<y\le m+q\), the original first anchor is \(y-3q-m+1\) and the
receiving index is \[
j_{\rm tail}=q-1-\ell.
\] Its anchor is \(y-m-2-3\ell\), exactly the common source anchor.
Equivalently, \[
j_{\rm tail}=j_{\rm stable}+q-y,
\qquad j_{\rm stable}=j_{\rm tail}+y-q.
\] The displayed bounds prove both indices valid. At \(d\ge m+1\),
\(h\ge m+q\), so every common row is an actual prefix row with index
\(y-1-\ell\). The maps retain sector, orientation, anchor and all three
cells, and consequently intertwine the original incidence maps. They
include the smallest exterior parameter, not only an eventual range.

\subsection{3. An exact eighteen-cell
scaffold}\label{an-exact-eighteen-cell-scaffold}

Start from the completed turn-8 \(q=3k\) replacement \(B_{m,k}\) and its
base stones. Translate them by \((-2,1)\) and delete one additional
original right stone before that translation, namely \(R(m-2-k,-k)\).
Denote the surviving translated stone family by \(S\). It has
\(t_m-3k-1\) members. The additional stone has base indices
\((r,s)=(0,k)\); it is distinct from the preceding three corner chains
because \(m\ge3k+3\).

The following table defines an original source subfamily
\(A_0\subset A\). Every entry uses the literal labels \((p,y,\ell)\)
above. Both endpoints of each range are included.

\begin{longtable}[]{@{}
  >{\raggedright\arraybackslash}p{(\columnwidth - 4\tabcolsep) * \real{0.3333}}
  >{\raggedright\arraybackslash}p{(\columnwidth - 4\tabcolsep) * \real{0.3333}}
  >{\raggedright\arraybackslash}p{(\columnwidth - 4\tabcolsep) * \real{0.3333}}@{}}
\toprule\noalign{}
\begin{minipage}[b]{\linewidth}\raggedright
Sector \(p\)
\end{minipage} & \begin{minipage}[b]{\linewidth}\raggedright
Depth \(\ell\)
\end{minipage} & \begin{minipage}[b]{\linewidth}\raggedright
Rows \(y\)
\end{minipage} \\
\midrule\noalign{}
\endhead
\bottomrule\noalign{}
\endlastfoot
0 & 0 & \(1,\ldots,m+2k-1\) \\
0 & \(1,\ldots,k-1\) & \(\ell+1,\ldots,m+2k-3\) and
\(m+2k-2+\ell,m+2k-1+\ell\) \\
0 & \(k\) & \(k+1,\ldots,m+3k-1\) \\
0 & \(k+1,\ldots,2k-1\) & \(\ell+1,\ldots,2k+1\) \\
1 & \(0,\ldots,k-3\) & \(\ell+1,\ldots,m+2k-6\) and
\(m+2k-3+\ell,m+2k-2+\ell\) \\
1 & \(k-2\) & \(k-1,\ldots,m+3k-4\) \\
1 & \(k-1\) & \(k,\ldots,2k-2\) and \(3k-1,3k,3k+1\) \\
1 & \(k,\ldots,2k-3\) & \(\ell+1,\ldots,2k-2\) \\
2 & \(0,\ldots,k-2\) & \(\ell+1,\ldots,m+2k-3\) and
\(m+2k-1+\ell,m+2k+\ell\) \\
2 & \(k-1\) & \(k,\ldots,m+3k-1\) \\
2 & \(k,\ldots,2k-2\) & \(\ell+1,\ldots,2k+1\) \\
2 & \(2k-1\) & \(2k\) \\
\end{longtable}

The depth ranges in each sector are disjoint except for the two
explicitly separated row intervals. Their separation follows
respectively from the strict gaps between \(m+2k-3\) and
\(m+2k-2+\ell\), between \(m+2k-6\) and \(m+2k-3+\ell\), and between
\(m+2k-3\) and \(m+2k-1+\ell\). The two ranges at sector 1, depth
\(k-1\), satisfy \(2k-2<3k-1\). All labels obey \(1\le y\le m+q\) and
\(0\le\ell<\min(y,q)\). Thus no old generator is repeated.

Define \[
\alpha(x,y)=\{(x,y),(x,y+1),(x+1,y-1)\}.
\] The eighteen residual cells \(Z\) are the following disjoint sets:

\begin{itemize}
\tightlist
\item
  \(\alpha(2k-2,4-4k-m)\);
\item
  \((3k-1,1-m),(3k-1,2-m),(3k,2-m),(3k,3-m),(3k+1,3-m),(3k+1,4-m)\);
\item
  the original right-stone cell set \(R(m-k-4,1-k)\);
\item
  the vertical triple \(V(m+2k-5,2k-4)\);
\item
  \((m+2k-2,2k-1),(m+2k-2,2k),(m+2k-1,2k)\).
\end{itemize}

The exact coefficient identity is \[
\partial\bigl((B_{m,k}+(-2,1))+S\bigr)+\mathbf1_Z
 =\mathbf1_{\Omega_\Delta}+\partial A_0.\tag{1}
\] Here \(B+(-2,1)\) means translation of every original tile, not
addition of a scalar. \path|source_table.py| writes this identity in six
independent Laurent variables. Its integer numerator cancels before
either parameter is specialized; Section 7 describes the complete
certificate mechanism.

The right side of (1) is the indicator of a disjoint union of original
cells. Every term on the left has nonnegative coefficient. Equation (1)
therefore proves the complete positive scaffold, including distinctness
of the residual cells. Adding the actual frozen family
\(A\setminus A_0\) gives a packing of
\(\Omega_\Delta\cup\operatorname{cells}(A)\) whose uncovered set is
exactly \(Z\).

\subsection{4. Original local homotopies and positive patch
identities}\label{original-local-homotopies-and-positive-patch-identities}

In addition to \(\alpha\), put \[
\beta(x,y)=\{(x,y),(x+1,y),(x+1,y+1)\}.
\] The following table states equalities of original cell partitions.
The old tiles together with the old holes partition exactly the same
support as the new tiles together with the new holes.

The \textbf{D move} replaces \(V(x+1,y),V(x+2,y-2)\) by
\(D(x,y),D(x,y+1)\) and sends \(\alpha(x,y)\) to \(\alpha(x+1,y+1)\).

The \textbf{U move} replaces \(V(x,y+2),V(x+1,y)\) by
\(V(x,y),V(x+1,y-1)\) and sends \(\alpha(x,y)\) to \(\alpha(x,y+3)\).

The \textbf{R move} replaces \(H(x+2,y),H(x+1,y+1)\) by
\(H(x,y),H(x,y+1)\) and sends \(R(x,y)\) to \(R(x+3,y)\).

The \textbf{B move} replaces \(V(x,y+1),V(x+1,y+2)\) by
\(V(x,y),V(x+1,y)\) and sends \(\beta(x,y)\) to \(\beta(x,y+3)\).

Each move is checked by expanding its two original three-cell
generators. If \(\kappa\) is new minus old, then \(\partial\kappa\) is
old-hole indicator minus new-hole indicator. Reversing the signed vector
and exchanging the two displayed positive partitions gives the inverse
local map. No division, averaged orbit or relaxed coefficient is used.

\subsubsection{4.1 The original intercorner channel, including zero
length}\label{the-original-intercorner-channel-including-zero-length}

Set \(L=m-3k-3\ge0\). Release \[
\begin{aligned}
O_L={}&\{V(3k+2+j,j-m):0\le j\le L\}\\
 &\cup\{V(3k+2+j,j+3-m):0\le j<L\},
\end{aligned}
\] and insert \[
N_L=\{D(3k+1+i,3-m+i+b):0\le i<L,\ b=0,1\}.
\] Their cell identity replaces the two holes \((3k+1,3-m),(3k+1,4-m)\)
by \((3k+2,-m),(3k+2,1-m)\) and \(\alpha(m-2,-3k)\). There are \(2L+1\)
released bones and \(2L\) inserted bones.

At \(L=0\), one original vertical bone is still released. The two old
holes are also two of the written new holes and cancel in the incidence
equation. The remaining three new holes are precisely that released
bone. Omitting this generator would incorrectly exclude the smallest
exterior parameter. The formula above retains it.

\subsubsection{4.2 Completion of the first nine-cell
endpoint}\label{completion-of-the-first-nine-cell-endpoint}

Apply \(k-1\) D moves beginning at \[
(x_i,y_i)=(2k-2+i,4-4k-m+i),\quad 0\le i<k-1.
\] Then apply \(k-1\) U moves beginning at \[
(3k-3,3-3k-m+3j),\quad0\le j<k-1.
\] The moving vacancy reaches \(\alpha(x,y)\) with \((x,y)=(3k-3,-m)\).
Together with it are the six fixed holes \[
(x+2,y+1),(x+2,y+2),(x+3,y+2),(x+3,y+3),(x+5,y),(x+5,y+1).
\] Release \(V(x+1,y),V(x+2,y-2),V(x+3,y-1),V(x+4,y)\). Insert \[
\begin{gathered}
D(x,y),D(x,y+1),D(x+1,y+1),D(x+3,y+3),\\
H(x+1,y+2),H(x+2,y+1),H(x+3,y).
\end{gathered}
\] The old four bones plus these nine holes and the new seven bones have
the same cell indicator. Expanding the offsets proves the identity. The
complete first-endpoint packet releases \(4k\) original vertical bones
and inserts \(4k+3\) bones.

\subsubsection{4.3 Transport and completion at the other
endpoint}\label{transport-and-completion-at-the-other-endpoint}

Move \(\alpha(m-2,-3k)\) by \(2k-2\) D moves. Move the right-stone
vacancy \(R(m-k-4,1-k)\) by \(k-1\) R moves. With \[
c=m+2k-4,\qquad b=1-k,
\] the lower six vacancies are now \[
\alpha(c,b-3)\cup R(c-3,b).
\] The following constant replacement, in coordinates relative to
\((c,b)\), replaces those six holes by \(\beta(-1,6)\).

Release the following fourteen original bones: \[
\begin{gathered}
H(-1,0),\ H(2,3),\\
V(-2,1),\ V(-1,1),\ V(-1,4),\ V(0,2),\ V(0,5),\\
V(1,-3),\ V(1,3),\ V(2,-5),\ V(2,-2),\\
V(3,-4),\ V(4,-3),\ V(4,0).
\end{gathered}
\] Insert the following fifteen original bones: \[
\begin{gathered}
D(0,-3),\ D(0,-2),\ D(1,-2),\ D(1,-1),\ D(2,-1),\\
H(-3,0),\ H(-3,1),\ H(-2,2),\ H(-2,3),\\
H(-1,4),\ H(-1,5),\ H(0,0),\ H(1,3),\\
V(4,-2),\ V(4,1).
\end{gathered}
\]

Thus fourteen old bones plus six holes equal fifteen new bones plus
three holes. The full relative translation and its inverse are
\((u,v)\mapsto(c+u,b+v)\) and \((x,y)\mapsto(x-c,y-b)\).

Apply \(k-4\) B moves at \[
(c-1,b+6+3j),\quad0\le j<k-4.
\] For \(k=4\) this is the empty sequence, not a negative-length family.
At the upper endpoint put \(b'=2k-5\). Relative to \((c,b')\), the nine
holes are \[
\beta(-1,0)\cup V(-1,1)\cup\{(2,4),(2,5),(3,5)\}.
\] Release \(H(0,2),H(0,3),V(3,2)\) and insert \[
H(-1,3),H(1,2),V(-1,0),V(0,0),V(2,3),V(3,3).
\] The corresponding eighteen-cell identity completes the patch.

The complete edit sequence releases \(2m+6k-2\) old scaffold bones and
inserts \(2m+6k+4\) bones. Its aggregate incidence difference is exactly
\(\mathbf1_Z\). All edit formulas, including the empty channel and empty
B-move range, appear in the shared literal \texttt{edits.json}.

\subsection{5. Why the original sources of every move
exist}\label{why-the-original-sources-of-every-move-exist}

A formal vacancy identity alone is insufficient to establish that a
proposed edit can be made in this benzel. The source maps are therefore
part of this proof.

Every released generator is identified either with an actual common
source label outside \(A_0\), or with a specified generator of the
translated turn-8 corner table. These inverse index formulas are stored
beside each old tile in \texttt{edits.json}. For example, the first two
D-move source families at the first endpoint have common labels \[
(p,y,\ell)=(1,2k-1+i,2k-2),\quad(1,2k+i,2k-1),
\quad0\le i<k-1.
\] Their original tiles are exactly the two displayed vertical bones,
because \[
F\rho H(x,y)=V(y,x).
\] The channel uses labels \((1,3k+2+j,k)\) and \((1,3k+2+j,k-1)\) in
the two stated ranges. The latter range stops one row earlier. This
retains the different endpoint roles.

The two sources of each R move at the other endpoint are original corner
families 6 and 7 (zero-based indices in the retained parent
\path|families.json|), sector 2, with respective indices \[
(i_6,j_6)=(k-1,k-1-n),\qquad i_7=k-2-n,
\quad0\le n<k-1.
\] The fixed lower patch uses the complementary family-6 index
\(j_6=0\). Its other exceptional corner indices and both upper-patch
indices are printed explicitly in the same table. All remaining sources
are common sector-1 vertical bones.

\path|ownership.py| substitutes every inverse index into the actual
original source cell formulas and checks equality of the three affine
cell coordinates. This is an identity in \((m,k,n)\), not a test at one
parameter. It also checks the domain of every inverse index and proves
that every common source used is outside \(A_0\).

The domain proof uses exact certificates. An inequality is represented
as an integer affine vector \(a\) with \(a\cdot(m,k,i,j,1)\ge0\). For a
required bound \(g\ge0\), the receipt expresses \(g\) as a nonnegative
rational combination of the exact domain inequalities and \(1\ge0\). For
an alleged intersection with \(A_0\) or a repeated original generator,
the receipt combines the stated domain and equality constraints into
\(0\ge1\). The verifier reconstructs the constraints and checks every
coefficient exactly using rational arithmetic. A numerical optimizer was
used only to discover multipliers; its status is not a proof
certificate.

For distinct atomic source families, both original row and depth, or
both parent corner indices, are included in this test. Within one
indexed family, a nonzero integer index coefficient proves injectivity.
The only shared parent family between the R-move source and the fixed
lower source has disjoint ranges \(j_6\ge1\) and \(j_6=0\).

Consequently every old generator occurs exactly once and is present in
the initial scaffold. Subtracting them leaves a nonnegative integral
tile vector. Inserting the listed positive new generators produces a
nonnegative vector. Section 7 proves its exact original cell boundary;
coefficient one at each target cell then proves a tiling. The proof does
not require an unconstructed positive lift.

\subsection{6. Return to the original full
benzel}\label{return-to-the-original-full-benzel}

Let \(B'\) be the final bone vector after all the edits. Combining their
incidence identity with (1) and the unchanged old common bones gives \[
\partial(B'+S)=\mathbf1_{\Omega_\Delta}+\partial A.\tag{2}
\] The source maps in Section 2 place every member of \(A\) in
\(\Gamma_\Delta(d,h)\). Thus \[
T(d,m,k)=\bigl(\Gamma_\Delta(d,h)\setminus A\bigr)\cup B'\cup S
\] is an original full tiling. The old frozen complement is disjoint
from the patch. Equation (2) proves the claimed counts: the patch has
\(\Delta\) stones, and its number of bones equals \(|A|\). The full bone
count is therefore unchanged from the parent packing, namely
\(3h(h+d)\).

This is a direct positive construction on the full displayed parameter
domain. No conditional theorem asserting that a future endpoint solution
would suffice is used.

\subsection{7. Exact parameter-preserving coefficient
proof}\label{exact-parameter-preserving-coefficient-proof}

The original cell-module isomorphism is \([(x,y)]\mapsto X^xY^y\) in
\(\mathbb Z[X^{\pm1},Y^{\pm1}]\), with coefficient extraction as
inverse. The original tile polynomials are \[
R=1+X+Y,\quad H=1+X+X^2,\quad V=1+Y+Y^2,\quad D=1+XY^{-1}+X^2Y^{-2}.
\] Both parameters are retained first in the six-variable Laurent ring
\[
\mathbb Z[X^{\pm1},Y^{\pm1},M_x^{\pm1},M_y^{\pm1},K_x^{\pm1},K_y^{\pm1}].
\] Specialization sends \(M_x,M_y,K_x,K_y\) to \(X^m,Y^m,X^k,Y^k\). Its
kernel consists of the four corresponding parameter relations. The
inverse of the induced quotient isomorphism is the inclusion of \(X,Y\).

An original indexed segment whose step is \((a,b)\) has its finite
geometric numerator divided by \(1-X^aY^b\). The source-table branch at
\(y=q\) is split explicitly: the prefix uses \(r_y=m-1\) and the tail
uses \(r_y=m\). No source mass, endpoint or multiplicity is inferred
from an eventual regime.

Two identities are evaluated independently: (1), and the aggregate edit
boundary equal to \(\mathbf1_Z\). \path|symbolic_certificate.py| expands
their integer Laurent numerators using the exact shared source and edit
tables. Every coefficient is zero before parameter specialization. All
denominators use original \(X,Y\) only and are nonzero. Localization is
injective because the Laurent ring is an integral domain; specialization
therefore returns the exact finite cell equations. All index lengths are
nonnegative on the parameter domain, with length zero retained
literally.

Together with Section 5, these evaluated identities supply the complete
positive partition proof. The finite tiling replay is an additional
check, not the source of the unbounded quantifiers.

\subsection{8. Actual Split-Zero complexes and
homotopies}\label{actual-split-zero-complexes-and-homotopies}

For the fixed \((d,m,k)\), index released original source subsets by
\(E\subseteq A\), with a separate bottom for empty physical support.
Define \[
\Lambda_E=\Omega_\Delta\cup\operatorname{cells}(E),\qquad
C_E^0=\mathbb Z[\text{original contained tiles}],\quad C_E^1=\mathbb Z[\Lambda_E],\quad C_E^2=0.
\] The maps are the original incidence and generator inclusions. The
reconstruction over \(G(\mathbb Z)\) retains \[
e\cdot(E,x)=(E,0),\qquad \tau\cdot(E,x)=(\bot,0).
\] Its differential square is the supported-zero map in the actual
target degree. The internal quotient coequalizes the original boundary
and its supported-zero companion. The proof of descent uses equality
modulo the original boundary submodule and absorption at the original
fibre zero, exactly as in the retained D1--D8 source.

For old and new bone boundaries
\(B_A,B_{B'}\subset\mathbb Z[\Lambda_A]\), the original representative
map gives \[
\ker\bigl(\mathbb Z[\Lambda_A]/B_A\to\mathbb Z[\Lambda_A]/(B_A+B_{B'})\bigr)
\cong B_{B'}/(B_A\cap B_{B'}).
\] Its inverse sends a killed representative \(a+b\) to \([b]\); two
decompositions differ precisely by the displayed intersection. Both
original bone families are matchings. The parent matching-graph
calculation consequently gives all primitive closed-component relations,
integral coordinate inverses and original-cell forest duals. No single
distinguished class is substituted for this full kernel.

The repaired representative is \[
z=\mathbf1_{\Omega_\Delta}-\partial S=\partial(B'-A).
\] Coefficient evaluation at the original cell \((m-2,0)\) annihilates
\(B_A\) and sends \(z\) to one. Indeed that cell has rank below the
deleted last shell and lies in \(\Omega_\Delta\), whereas its first
coordinate exceeds the maximum first coordinate of the translated
all-stone tiling. It is therefore a split integral class in the old
quotient before being killed by the displayed original new boundaries.
Its image is its own supported zero, not absence.

The local D, U, R and B moves give explicit degree-minus-one cochain
homotopies: if \(f,g\) mark their old and new hole cells, then
\(\partial H=f-g\), where \(H\) is the actual new-minus-old bone vector.
Their sums telescope through both endpoints. On finite collections of
the actual path generators these identities commute with the inclusion
maps, and hence reconstruct over \(G(\mathbb Z)\) with the support
labels unchanged.

Finally, set \(v_E=\mathbf1_{\Lambda_E}\) and \[
\widehat d_E(r,n)=\partial n-rv_E.
\] For \(E\subseteq E'\), let \(\eta_{EE'}\) be the newly released
original bones. Then \[
\widehat j_{EE'}(r,n)=(r,jn+r\eta_{EE'}),\qquad
v_{E'}=jv_E+\partial\eta_{EE'}.
\] Expanding proves the cochain identity and composition. Charge-one
nonnegative integral cycles are exactly the positive tiling fibre.
Adjoining the frozen outer source bones gives the linear augmented map
into the full benzel complex; deleting those same frozen bones is its
inverse on the specified containing-tiling fibre. Thus the signed
homotopies, actual support comparisons and positive completion are
related by displayed morphisms.

\subsection{9. Scope, provenance and
continuation}\label{scope-provenance-and-continuation}

The parent mathematical input is the preserved turn-8 source at Commons
commit \path|20cb022dcfda7f37b3adec13b5619be8f13be3e8|, and its complete
stable-packing source. The Split-Zero source remains
\path|zeta-function-research-reader@1c8ec52c85c173adc9f8403a8a26914955e2f5a9|,
especially \path|support_diagrams.tex| D1--D8 and
\path|SplitZeroComplex.lean|. These are dependencies, not claimed new
constructions of general homological algebra.

This continuation supplies \(q=3k+1\) for \(k\ge4\), all \(m\ge3k+3\)
and all \(d\ge m\). It does not supply the unbounded \(q=3k+2\) positive
endpoint, nor does it replace the earlier separate proofs for small
deletion values. Those remain the next original positive-source
calculations. No independent acceptance, Lean certificate or full P7
resolution is asserted.

The archive retains exact integer coefficient and affine source
certificates, original-coordinate code, \path|normal/optimized| replay
receipts, and the unchanged earlier source editions. Discovery searches
and their numerical statuses are not dependencies of the final
constructor.

\end{document}
